% SHARD-ID: AQM-72-QSA-BOSON-FERMION-LAYER
% SHARD-TITLE: Bosonic and Fermionic Association Layer
% SHARD-SUMMARY: Reviews the finite bosonic association via real symplectic map spaces and CCR/Weyl representations.
% SHARD-SUMMARY: Reviews the finite fermionic CAR layer and the Grassmann-Weyl displacement calculus with active local source locators.
% SHARD-SUMMARY: Records parafermion, fusion-category, and finite AND/OR comparison variants as gaps rather than results.
% SHARD-KEYWORDS: quantum-system association, bosons, fermions, CCR, CAR, Grassmann-Weyl, parafermion gap, fusion-category gap

\section{Bosonic and Fermionic Association Layer}
\label{sec:qsa-boson-fermion-layer}

\subsection{Status and Local Anchors}

This is QSA Step 08 from
\path{docs/quantum_system_associations/PLAN.md}.  It is a review shard.  It
adds no new examples, no new source claims, no new run artifacts, and no
comparison theorem with the finite AND/OR grammar of AQM-70.

The vocabulary and evidence discipline are those of
\texttt{CONVENTIONS.md} convention \texttt{(ap)} and AQM-65.  The canonical
boson/fermion comparison convention is \texttt{CONVENTIONS.md} convention
\texttt{(s)}.  The main report anchors are
AQM-19 for finite CCR/CAR field-label tests, AQM-20 for the abstract
Grassmann displacement ray algebra and Cahill--Glauber representation, and
AQM-21 for the unitarity, displacement, ray-product, and ordinary
Hilbert-space boundary derivations.

The active source manifest is
\path{references/canonical_fields/SOURCES.md}.  For the bosonic CCR layer it
registers Derezinski's
\path{references/canonical_fields/derezinski_math_ph_0511030_source/derezinski.tex}
at lines 1078--1135 for Weyl-form CCR, 1197--1284 for the Schrodinger
representation and finite-dimensional Stone--von Neumann statement, and
2251--2301 for bosonic Fock CCR fields; it also registers Bekka's
\path{references/canonical_fields/bekka_2502_00387_source/CCR-GeneralRings-v5.tex}
at lines 220--231, 253--273, 287--350, and 603--633 for the finite real or
local-field CCR uniqueness layer.  For the fermionic CAR layer it registers
Derezinski's
\path{references/canonical_fields/derezinski_math_ph_0511030_source/derezinski.tex}
at lines 1437--1524 for CAR definitions,
1612--1696 for finite-dimensional CAR uniqueness, 1803--1964 and
2098--2131 for Fock and creation/annihilation operators, and 2629--2662 for
fermionic Fock CAR fields.

For the Grassmann displacement layer, the active Cahill--Glauber source is the
official arXiv e-print extraction
\path{references/canonical_fields/cahill_glauber_physics_9808029v1_source/fxxx.tex}.
The active locators are lines 323--335 for CAR and vacuum, 337--365 for
Grassmann variables and adjunction, 603--643 for nilpotency and parity of
Grassmann functions, 367--390 for \(D(\gamma)\), 391--414 for displacement
identities, 416--464 for the ray-product formula, 815--833 for completeness,
and 835--854 for the caveat that physical quantum operators do not themselves
involve Grassmann variables.  The historical APS/marker chain recorded in the
manifest is source-integrity provenance only and is not used as active
evidence here.

The same canonical-field manifest registers Keyl--Schlingemann's
\path{references/canonical_fields/keyl_schlingemann_0906_2929_source/fidelity_arXiv.tex}
at lines 253--260, 270--355, 441--492, and 496--531 for GAR as a twisted
Grassmann/CAR framework and for the caveat that the Grassmann phase-space
right hand side is not an ordinary complex scalar.  The supergeometry manifest
\path{references/supergeometry/SOURCES.md} registers Witten's
\path{references/supergeometry/witten_1209_2199_source/supernotes.tex} at
lines 652--660 for odd variables quantized by Clifford anticommutation
relations and lines 671--704 for the \(\Pi TM\) parity-reversal and
Clifford-action dictionary.  This shard uses those supergeometry locators only
as background anchors for the parity/Clifford boundary; it does not import a
new finite-field supergeometry example.

\subsection{Catalogue Slots}

In QSA language, the locally available boson/fermion layers are the following
review entries.

\begin{center}
\small
\begin{tabular}{p{0.22\linewidth}p{0.31\linewidth}p{0.28\linewidth}p{0.13\linewidth}}
\toprule
Layer & Added input data & Kinematical output & Status \\
\midrule
Finite bosonic CCR field &
Finite site set \(X\), real symplectic target
\(V_B=\mathbb R_q\oplus\mathbb R_p\), and
\(E_B=\operatorname{Map}(X,V_B)\) with the summed pointwise form. &
Regular irreducible Weyl/CCR representation, Schrodinger model
\(L^2(\mathbb R^{|X|})\), and equivalent finite-mode bosonic Fock model. &
reviewed from AQM-19 T1 \\
\addlinespace
Finite fermionic CAR field &
Finite site set \(X\), finite-dimensional complex internal Hilbert space
\(K\), \(Z=\operatorname{Map}(X,K)\), and
\(Y_F=\operatorname{Re}(Z\oplus\overline Z)\) with positive scalar product. &
CAR/Clifford field operators on the antisymmetric Fock space
\(\Gamma_{\mathrm a}(Z)\). &
reviewed from AQM-19 T2 \\
\addlinespace
Grassmann-Weyl displacement calculus &
Finite CAR mode set \(I\), Grassmann coefficient algebra, and the graded
Grassmann-extended CAR algebra. &
Unitary elements \(D(\gamma)\) in the extended algebra with
Grassmann-valued ray product and CAR displacement identities. &
reviewed from AQM-20/21 \\
\bottomrule
\end{tabular}
\end{center}

The entries are not bare finite-set associations.  Each begins only after the
displayed extra data have been supplied.

\subsection{Finite Bosonic Association}

AQM-19 D1--D4 and convention \texttt{(s)} place a finite bosonic field in the
same broad field-label slot as AQM-14, but with real symplectic labels rather
than finite-field labels.  The label space is
\[
  E_B=\operatorname{Map}(X,\mathbb R_q\oplus\mathbb R_p),
\]
and the form is the finite sum
\[
  \Omega_B((q,p),(q',p'))
  =
  \sum_{x\in X}(p_x q'_x-p'_x q_x).
\]
AQM-19 L1 checks that this is a finite-dimensional real symplectic space.
AQM-19 L2 then invokes the locally registered Stone--von Neumann sources to
identify an irreducible regular Weyl-form representation with the Schrodinger
representation on \(L^2(\mathbb R^{|X|})\), and AQM-19 T1 records the finite
bosonic field comparison.  This is a sourced finite kinematical CCR
association, not a dynamics statement.

\subsection{Finite Fermionic CAR Association}

AQM-19 D5--D7 use different input data for standard finite fermions.  After
choosing a finite-dimensional complex internal one-particle Hilbert space
\(K\), set
\[
  Z=\operatorname{Map}(X,K),
  \qquad
  Y_F=\operatorname{Re}(Z\oplus\overline Z).
\]
The scalar product is
\[
  \alpha((z,\bar z),(w,\bar w))=\operatorname{Re}\langle z,w\rangle_Z.
\]
On the antisymmetric Fock space
\[
  \Gamma_{\mathrm a}(Z)=\bigoplus_k \wedge^k Z,
\]
AQM-19 defines \(\phi(z,\bar z)=a^*(z)+a(z)\), proves the exterior-algebra
CAR in L3, and proves
\[
  \{\phi(y),\phi(y')\}=2\alpha(y,y')I
\]
in L4.  AQM-19 T2 therefore records the finite standard fermionic canonical
field as a CAR/Clifford association.  It is not obtained by taking the
bosonic real symplectic map-space recipe and replacing the target by
Grassmann variables.

\subsection{Grassmann-Weyl Displacement Review}

The Grassmann layer is available only inside a Grassmann-extended CAR algebra.
AQM-20 defines a finite CAR algebra generated by \(a_i,a_i^\dagger\), a
Grassmann coefficient algebra generated by \(\gamma_i,\gamma_i^*\), and the
graded product algebra in which the Grassmann generators anticommute with the
odd CAR generators.  It then defines
\[
  K(\gamma)=\sum_i(a_i^\dagger\gamma_i-\gamma_i^*a_i),
  \qquad
  D(\gamma)=e^{K(\gamma)}.
\]
AQM-21 proves \(K(\gamma)^\dagger=-K(\gamma)\), so \(D(\gamma)\) is unitary
in the extended algebra; proves that conjugation displaces the CAR generators
by \(a_i\mapsto a_i+\gamma_i\) and
\(a_i^\dagger\mapsto a_i^\dagger+\gamma_i^*\); and proves the ray product
\[
  D(\alpha)D(\beta)
  =
  D(\alpha+\beta)
  \exp\!\left(
    \frac12\sum_i(\beta_i^*\alpha_i-\alpha_i^*\beta_i)
  \right).
\]
AQM-21 T2 is the boundary: the multiplier is an even Grassmann unit, not a
complex \(U(1)\) phase, so this is not an ordinary scalar-projective unitary
representation on a complex Hilbert space unless the Grassmann coefficient
algebra has first been included and tracked.

\subsection{Boundary With QSA Grammar}

AQM-70 supplies a finite AND/OR carrier grammar:
\(\operatorname{AND}=\otimes_{\mathbb C}\) and
\(\operatorname{OR}=\oplus\).  That grammar records finite vector-space
bookkeeping and explicitly defers symmetric, antisymmetric, para, Fock, and
fusion-category enhancements.  This shard does not identify the bosonic CCR
or fermionic CAR layers with that grammar.

The finite bosonic entry requires a real symplectic map space plus a regular
CCR/Weyl representation.  The finite fermionic entry requires a positive
CAR/Clifford label space and antisymmetric Fock representation.  The
Grassmann displacement entry requires the extended CAR coefficient algebra and
has a Grassmann-valued multiplier.  No map, quotient, invariant, or local
derivation currently compares these outputs with AQM-70 as an agreement or
equivalence result.

\subsection{Gaps}

The following variants are not results in the current catalogue.

\begin{itemize}
\item \textbf{Parafermions and para-statistics.}
No parafermion source locator, convention, mode algebra, permutation-sector
rule, or checked example is registered for this QSA step.  Convention
\texttt{(at)} already marks para-statistics as a separate enhancement, so any
parafermion association is a gap.
\item \textbf{Fusion-category and Levin-Wen/string-net variants.}
Convention \texttt{(h)} leaves the Levin-Wen fusion-category gauge,
pivotal/spherical structure, and \(F/R\)-data unfixed, and the QSA plan
assigns this to Step 09.  No fusion-category endpoint is asserted here.
\item \textbf{Finite AND/OR equivalence.}
AQM-70 and this shard have no recorded comparison map or invariant.  Any
claim that CCR/CAR/Fock layers agree with, reduce to, or are generated by the
finite AND/OR grammar remains unsupported.
\item \textbf{Arithmetic Grassmann target replacement.}
AQM-19 L5 rejects only the strict scalar-projective target-replacement
proposal: Grassmann variables do support fermionic phase-space calculus, but
inside GAR or Grassmann-extended CAR algebra rather than as an ordinary
finite-field or real symplectic Weyl target.  No new arithmetic Grassmann
association is introduced here.
\end{itemize}

\subsection{Conclusion}

QSA Step 08 places the existing canonical-field work into the association
catalogue as three reviewed layers: real symplectic CCR fields, CAR/Clifford
fermion fields, and Grassmann-Weyl displacement calculus in the extended CAR
algebra.  The unsupported variants remain explicit gaps, and no equivalence
with the elementary finite AND/OR grammar is claimed.
